% ============================================================
%  SECTION 3 (was A): The rupee's account book   (Member 1)
%  Defines "current account deficit" used in Sections 9 and 11
% ============================================================
\section{The rupee's account book: the balance of payments}\label{sec:bop}

Every dollar that enters or leaves India is written down in the \defterm{balance of payments (BoP)}{The record of all economic transactions between a country's residents and the rest of the world over a period.}, the record of all economic transactions between the residents of a country and the rest of the world over a period such as a year. The \defterm{current account}{The part of the BoP that records trade in goods and services, income and transfers.} records trade in goods and services, along with income and transfers such as remittances. The \defterm{capital account}{The part of the BoP that records the buying and selling of assets, such as FDI, FPI and loans.} records the buying and selling of assets, such as FDI, FPI and foreign loans. When a country spends more foreign currency on its current account than it earns, it runs a \defterm{current account deficit (CAD)}{The gap when a country's payments abroad for goods, services, income and transfers exceed its receipts.}. The gap must then be paid for by a net inflow on the capital account or by drawing down the country's reserves.\cite{ncert2007}

India imports far more goods than it exports, mainly oil, gold and electronics. But India also sells services such as software abroad, and Indians working abroad send money home. The dollars from these two sources cover most of the gap. In 2025-26 these and other \defterm{invisibles}{Trade in services, income and transfers, as opposed to trade in goods.} (services, income and transfers) brought in a net US\$312 billion, leaving a current account deficit of only US\$25.2 billion, about 0.6 per cent of GDP.\cite{rbi2026b} The capital account could not fill even this gap, because foreign portfolio investors took out a net US\$16.4 billion. So the RBI drew about US\$23.6 billion from its reserves \seesec{sec:reserves}.
